%%%PAGE 6 rot=0 legibility=good summary=康普顿效应；粒子性与波动性/实物波·概率波提纲及应用；巴尔末系；玻尔模型三条假设；氢原子能级图、半径与能级示意图；能级跃迁关系式
%%%BLOCK id=006-01 ch=17 sec=波粒二象性·康普顿效应 kind=knowledge alt=none cont=none y=0.02-0.10
\textbf{康普顿效应}\quad 波长$\lambda_0$的X射线照射石墨，散射成分中出现波长大于$\lambda_0$的成分。

光子与电子作用时，动量能量都守恒\quad $P=\dfrac{h}{\lambda}$
%%%END
%%%BLOCK id=006-02 ch=17 sec=波粒二象性·物质波 kind=summary alt=none cont=none y=0.10-0.41
\begin{itemize}
\item 粒子性
  \begin{itemize}
  \item 两个效应
    \begin{itemize}
    \item 光电效应
    \item 康普顿效应
    \end{itemize}
  \item 三个量
    \begin{itemize}
    \item 动量
    \item 能量
    \item 电量
    \end{itemize}
  \end{itemize}
\end{itemize}

\begin{itemize}
\item 粒子性：少，$\lambda$小，与物质作用
\item 波动性：多，$\lambda$大，光传播过程
\end{itemize}

\begin{itemize}
\item 粒子波动性
  \begin{itemize}
  \item 实物波（德布罗意波）
  \item 概率波（微观粒子）
  \end{itemize}
\item 波动性
  \begin{itemize}
  \item $\lambda=\dfrac{h}{p}$\quad $\nu=\dfrac{\varepsilon}{h}$\qquad 2个量
  \item 干涉\quad 衍射\quad 偏振\qquad 3个现象
  \end{itemize}
\end{itemize}

电子动量\quad $P=\sqrt{2mUe}$ % 原稿位于页面右上侧（波动性大括号上方）；首字“电”有涂改叠写

\begin{itemize}
\item 应用
  \begin{itemize}
  \item 电子显微镜（电子$m$大\ $P$大\ $\lambda$小\ 难衍射） % 原稿“难衍射”写在“λ小”下方一行，其后接右括号
  \item 热慢中子衍射 % 原稿“热”与“慢”之间被大括号的竖线隔开
  \item 提高精度：质子显微镜，增大速度
  \end{itemize}
\end{itemize}
%%%END
%%%BLOCK id=006-03 ch=17 sec=氢原子光谱 kind=formula alt=none cont=none y=0.43-0.46
\textbf{巴尔末系}\quad $\dfrac{1}{\lambda}=R\left(\dfrac{1}{2^2}-\dfrac{1}{n^2}\right)$
%%%END
%%%BLOCK id=006-04 ch=17 sec=玻尔模型 kind=knowledge alt=none cont=none y=0.46-0.63
\textbf{玻尔能级模型}\quad 三条假设：
\begin{itemize}
\item[①] \textsuperscript{原子}轨道/能量（能级）\textsuperscript{分立的}不连续\quad 电子在\textsuperscript{加速}高速运动但不向外辐射能量（电磁波）\quad 称为定态 % “加速”二字小字写在“高速”上方，疑为对“高速”的订正/补充
\item[②] 原子从一定态\textsuperscript{跃迁}到另一定态\quad 吸收或辐射\unclear{一定}的电磁波能量\quad 由两个能级差决定
\item[③] 原子能级$=$电子动能$+$系统电势能
\end{itemize}

（原子跃迁$=$电子轨道的跃迁） % 原稿“原子”与“跃迁”之间有一小块涂改液涂白

高轨低速大周期、大机大势\textsubscript{电势能}\ 大能量\textsubscript{能级}
%%%END
%%%BLOCK id=006-05 ch=17 sec=氢原子能级 kind=knowledge alt=none cont=none y=0.64-0.88
\textbf{氢原子能级图}

\[
r_n=n^2r_1,\qquad E_n=\frac{1}{n^2}E_1\quad(E_1=-13.6\unit{eV})
\]

%FIG p006_a 0.10 0.738 0.46 0.885
\notefig[0.35]{figures/p006_a.png}

氢原子只发出4种可见光\ $\to$\ $6\to2$\quad $5\to2$\quad $4\to2$\quad $3\to2$ % 原稿一条长箭头从“只发出4种可见光”指向页面右上方的“6→2  5→2  4→2  3→2”

%FIG p006_b 0.47 0.669 0.63 0.777
\notefig[0.25]{figures/p006_b.png}

间隙越来越大 % 图 p006_b（半径示意图）内亦含此字

%FIG p006_c 0.655 0.652 0.97 0.795
\notefig[0.45]{figures/p006_c.png}

间距越来越小 % 图 p006_c（能级示意图）内亦含此字

$n=1$叫基态\quad $n>1$叫激发态

$n=2$\ 第一激发态…
%%%END
%%%BLOCK id=006-06 ch=17 sec=氢原子能级 kind=formula alt=none cont=none y=0.88-1.00
能级$\ge0$叫电离\qquad 能级间距$=$代表吸收或辐射电磁波/光子

\[
h\nu_1=h\nu_2+h\nu_3\qquad \nu_1=\nu_2+\nu_3 % CHECK: 右式等号手写形似“≠”，按物理关系应为“=”，照原样按“=”转录
\]
\[
h\frac{c}{\lambda_1}=h\frac{c}{\lambda_2}+h\frac{c}{\lambda_3}\qquad \frac{1}{\lambda_1}=\frac{1}{\lambda_2}+\frac{1}{\lambda_3}
\]
%%%END
%%%PAGEEND 6
